\documentclass[preprint,12pt]{elsarticle}
\usepackage{textcomp}
\usepackage{xcolor}
\usepackage{graphicx}
\usepackage{booktabs}
\usepackage{amsmath}
\usepackage{hyperref}
\usepackage{microtype}

% ============================================================
% results.tex -- single source of truth for every number.
% Change a value here, it updates everywhere in the paper.
% Never type a number directly into the prose.
% ============================================================

% ---------- Dataset and pipeline parameters ----------
\newcommand{\nSubjects}{43}
\newcommand{\nSessions}{3}
\newcommand{\sessionDays}{1{,} 8{,} and 29}
\newcommand{\fsClinical}{2048}
\newcommand{\fsCommodity}{1000}
\newcommand{\nWindows}{37}
\newcommand{\nRows}{100{,}233}
\newcommand{\windowMs}{200}
\newcommand{\stepMs}{100}
\newcommand{\transientMs}{200}
\newcommand{\zcEpsilon}{20}          % microvolts
\newcommand{\hpCutoff}{20}           % Hz

% ---------- Channel derivation evidence ----------
\newcommand{\rawCorr}{0.987}
\newcommand{\bipolarCorr}{0.139}
\newcommand{\interRingCm}{2}

% ---------- Quantization ranges (mV) ----------
\newcommand{\vmaxTwoCh}{0.3329}
\newcommand{\vmaxEightCh}{0.6407}
\newcommand{\vmaxEightChResampled}{0.6395}
\newcommand{\lsbCommodity}{0.651}    % microvolts
\newcommand{\lsbEightCh}{1.252}      % microvolts

% ---------- LDA 2x2 (headline table) ----------
\newcommand{\ldaClinWithin}{96.2\%}
\newcommand{\ldaCommWithin}{86.1\%}
\newcommand{\ldaClinCross}{82.5\%}
\newcommand{\ldaCommCross}{71.4\%}
\newcommand{\ldaClinDrop}{13.7}
\newcommand{\ldaCommDrop}{14.7}
\newcommand{\ldaInteractionP}{0.59}
\newcommand{\ldaInteractionCI}{[-5.3{,} +7.2]}

% ---------- Random Forest ----------
\newcommand{\rfClinDayOne}{97.2\%}
\newcommand{\rfClinDayEight}{91.5\%}
\newcommand{\rfClinDayTwentyNine}{91.8\%}
\newcommand{\rfCommDayOne}{86.9\%}
\newcommand{\rfCommDayEight}{72.1\%}
\newcommand{\rfCommDayTwentyNine}{73.0\%}
\newcommand{\rfClinDrop}{5.4}
\newcommand{\rfCommDrop}{14.0}
\newcommand{\rfInteractionP}{0.0091}
\newcommand{\rfInteractionCI}{[+2.8{,} +14.5]}
\newcommand{\rfInteractionNormP}{0.0049}
\newcommand{\bonferroniAlpha}{0.01}

% ---------- Model-ceiling spread (post-hoc) ----------
\newcommand{\clinModelSpread}{9.3}
\newcommand{\commModelSpread}{2.3}
\newcommand{\deployCost}{1.6}

% ---------- Ablation: Day 1 difference from clinical ----------
\newcommand{\rateDayOneDiff}{+0.16}
\newcommand{\rateDayOneCI}{[-0.06{,} +0.58]}
\newcommand{\depthDayOneDiff}{-0.09}
\newcommand{\depthDayOneCI}{[-0.27{,} +0.00]}
\newcommand{\rateDepthDayOneDiff}{+0.18}
\newcommand{\rateDepthDayOneCI}{[-0.01{,} +0.72]}
\newcommand{\spatialDayOneDiff}{-11.10}
\newcommand{\spatialDayOneCI}{[-14.67{,} -8.31]}
\newcommand{\commodityDayOneDiff}{-10.23}
\newcommand{\commodityDayOneCI}{[-14.82{,} -7.25]}

% ---------- Ablation: cross-session drop difference from clinical ----------
\newcommand{\rateDropDiff}{-0.18}
\newcommand{\rateDropCI}{[-1.78{,} +0.48]}
\newcommand{\depthDropDiff}{-0.10}
\newcommand{\depthDropCI}{[-0.29{,} +0.03]}
\newcommand{\rateDepthDropDiff}{-0.27}
\newcommand{\rateDepthDropCI}{[-1.88{,} +0.34]}
\newcommand{\spatialDropDiff}{+9.09}
\newcommand{\spatialDropCI}{[+3.32{,} +14.66]}
\newcommand{\commodityDropDiff}{+8.57}
\newcommand{\commodityDropCI}{[+2.78{,} +14.38]}

% ---------- Commodity vs spatial-only (direct paired difference) ----------
\newcommand{\commVsSpatialDayOne}{+0.87}
\newcommand{\commVsSpatialDayOneCI}{[-0.67{,} +2.42]}
\newcommand{\commVsSpatialDrop}{-0.52}
\newcommand{\commVsSpatialDropCI}{[-3.06{,} +1.26]}

% ---------- SHAP ----------
\newcommand{\shapMigrationDayOne}{15.0\%}
\newcommand{\shapMigrationDayTwentyNine}{14.9\%}
\newcommand{\shapMigrationP}{0.92}
\newcommand{\shapExcessClinical}{0.0048}
\newcommand{\shapExcessSpatial}{0.0086}
\newcommand{\shapExcessCommodity}{0.0131}
\newcommand{\shapSimilarityClinical}{98.7\%}   % clinical, day 1 vs day 29
\newcommand{\shapSimilarityCommodity}{97.8\%}  % commodity, day 1 vs day 29

% ---------- Deployability (bytes) ----------
\newcommand{\flashLDA}{172}
\newcommand{\flashSVM}{172}
\newcommand{\flashKNN}{18{,}379}
\newcommand{\flashRFmin}{1{,}334}
\newcommand{\flashRFmax}{29{,}334}
\newcommand{\flashXGBmin}{2{,}172}
\newcommand{\flashXGBmax}{20{,}957}
\newcommand{\sramNaive}{1600}
\newcommand{\sramIntSixteen}{800}
\newcommand{\rfSizeCorr}{-0.456}
\newcommand{\rfSizeCorrP}{0.002}

% ---------- Confusion matrices and failure modes ----------
\newcommand{\cmClinGestureError}{16.6}       % percentage points
\newcommand{\cmClinRestError}{8.2}           % percentage points
\newcommand{\cmCommGestureError}{55.1}       % percentage points
\newcommand{\cmCommRestError}{26.1}          % percentage points
\newcommand{\cmCommCloseToExt}{30.1\%}
\newcommand{\cmCommExtToClose}{25.0\%}
\newcommand{\cmClinRestDrop}{-1.0}
\newcommand{\cmCommRestDrop}{-13.5}
\newcommand{\cmRestMisclass}{17.2\%}
\newcommand{\fOneFNineVmax}{0.1622}          % mV (from calibration)
\newcommand{\fSixFFourteenVmax}{0.9616}      % mV (from calibration)

% ---------- Embedded deployability constraints ----------
\newcommand{\measuredFlash}{8{,}344}
\newcommand{\flashBudget}{23{,}912}
% HISTORICAL, from the profiling sketch. Superseded by \verifySRAM{} (479,
% measured on the verification firmware). Do not use in the manuscript.
%\newcommand{\measuredSRAM}{648}
% HISTORICAL, superseded by \accumBytesMeasured{} (84, measured via sizeof).
% Retained for reference only. Do not use in the manuscript.
%\newcommand{\accumSRAM}{88}
\newcommand{\wcetTick}{1{,}612}
\newcommand{\latencyTotal}{202}

% ---------- Embedded implementation verification ----------
\newcommand{\nVerifyWindows}{666}
\newcommand{\nVerifyTrials}{18}
\newcommand{\nFilterSamples}{4{,}000}
\newcommand{\abStdErr}{1.9 \times 10^{-6}}
\newcommand{\abFiltErr}{8.9 \times 10^{-4}}
\newcommand{\abScoreErr}{3.2 \times 10^{-5}}
\newcommand{\abMargin}{0.0094}
\newcommand{\subBlockBytes}{21}
\newcommand{\accumBytesMeasured}{84}
\newcommand{\filterStateBytes}{32}
\newcommand{\dspStateBytes}{116}
\newcommand{\verifySRAM}{479}
\newcommand{\verifyFlash}{10{,}374}

\usepackage{fancyhdr}
% geometry is loaded only so fancyhdr has room for the header rule and
% distribution footer; elsarticle's default text block is otherwise preserved.
\usepackage[headheight=15pt, includeheadfoot]{geometry}

\journal{Biomedical Signal Processing and Control}
% --- Define the Single Footer Content ---
\newcommand{\distributionFooter}{%
  \begin{center}
    \small Approved for Public Release; Distribution Unlimited: PA CLEARANCE PENDING - DO NOT DISTRIBUTE. \\
    \vspace{2pt}
    \thepage
  \end{center}%
}

% --- Global Header & Footer Configuration ---
\pagestyle{fancy}
\fancyhf{} % Clear default header and footer fields
\renewcommand{\headrulewidth}{0.2pt}
\renewcommand{\footrulewidth}{0pt}

\lhead{\small \textit{Biomedical Signal Processing and Control}}
\cfoot{\distributionFooter}

% --- Override ALL Page Styles (plain, pprintTitle) to match 'fancy' ---
\fancypagestyle{plain}{\fancyhf{}\renewcommand{\headrulewidth}{0.2pt}\lhead{\small \textit{Biomedical Signal Processing and Control}}\cfoot{\distributionFooter}}

\makeatletter
\let\ps@pprintTitle\ps@fancy % Force elsarticle's title page style to be identical to fancy
\makeatother

\begin{document}

\begin{frontmatter}

\title{Do Hardware Constraints Compound Cross-Session Degradation? A Controlled Ablation of Channel Count, Sample Rate, and Quantization in Surface EMG Gesture Classification}

\author[a,b]{Adnan Hossain}
\author[b]{Jing Lin\corref{cor1}}
\cortext[cor1]{Corresponding author}

\affiliation[a]{organization={The City College of New York},
                city={New York},
                state={NY},
                country={USA}}

\affiliation[b]{organization={Air Force Research Laboratory},
                city={Rome},
                state={NY},
                country={USA}}


\thispagestyle{fancy} % Added
\begin{abstract}
Continuous surface electromyography (sEMG) gesture classification suffers from cross-session accuracy degradation due to electrode shift and signal non-stationarity. While commodity hardware makes edge deployment practical, it imposes strict compromises in spatial resolution, sampling rate, and signal quantization. It remains unclear which constraint drives cross-session failure, and whether these limits compound a classifier's fragility to temporal drift. Using the open-access GRABMyo dataset (\nSubjects{} subjects), we systematically degraded an 8-channel, floating-point clinical baseline into a 2-channel, 10-bit, \fsCommodity{}~Hz commodity configuration across a 28-day evaluation interval. Because all six conditions derive from identical source recordings, penalties were isolated through exactly paired comparisons. Spatial resolution is the dominant bottleneck: reducing the system from 8 to 2 channels lowered Day-1 accuracy by 11.10 percentage points and added 9.09 percentage points of cross-session degradation. Reducing sampling rate and bit depth altered cross-session degradation by less than 2 percentage points, with tightly bounded confidence intervals. The compounding effect of hardware constraints is model-dependent. Random Forest classifiers show the smallest cross-session loss on 8 channels, an advantage that collapses on 2 channels, compressing algorithmic variance and yielding diminishing returns for complex models. Feature attribution using SHapley Additive exPlanations (SHAP) showed a profile cosine similarity of 0.978 over 28 days despite a 14.0-percentage-point accuracy drop, consistent with shifted inputs evaluated against fixed decision boundaries. Profiling on an ATmega328P microcontroller shows that 2-channel deployment favors minimal architectures such as Linear Discriminant Analysis, shifting the bottleneck to spatial coverage and windowing latency.
\end{abstract}

\begin{keyword}
surface electromyography \sep gesture recognition \sep cross-session \sep
embedded systems \sep hardware constraints \sep feature attribution
\end{keyword}
\end{frontmatter}

\section{Introduction}
Surface electromyography (sEMG) records the electrical activity generated by contracting muscle through electrodes placed on the skin. Because the resulting signal carries information about intended movement, sEMG pattern recognition has become an established approach for controlling prosthetic limbs, rehabilitation devices, and wearable human--machine interfaces~\cite{ASGHARIOSKOEI2007275}. The classical pipeline---segment the signal into short windows, extract time-domain features, and classify with a conventional model---has been reproduced across many laboratories and remains computationally light enough for real-time operation~\cite{204774, 1206493}.

Two obstacles have nonetheless limited deployment outside controlled settings. First, classification accuracy degrades substantially when a user removes the device and puts it back on days later, driven by electrode displacement, skin-interface changes, and other sources of signal non-stationarity~\cite{10.3389/fnbot.2018.00058, 8009405}. Second, commodity acquisition hardware imposes constraints absent from research-grade amplifiers: fewer electrodes, lower sampling rates, and coarser analog-to-digital conversion. While each obstacle is well studied in isolation, they are rarely evaluated together. Longitudinal work generally uses multi-channel laboratory equipment, whereas hardware-constrained and embedded work evaluates performance within a single recording session~\cite{Pizzolato2017-zd, 7974234}.

This separation leaves an open question with direct engineering consequences. A designer building a low-cost sEMG interface must allocate a limited budget across electrode count, converter resolution, and sampling rate. Existing work treats commodity hardware as a single aggregate variable and therefore provides little guidance on which of these compromises actually costs accuracy---particularly once the device is used across multiple days without recalibration.

We address this by degrading a clinical baseline along three hardware axes independently and measuring the consequences over a 28-day interval. Using the GRABMyo dataset~\cite{pradhan2022grabmyo, PhysioNet-grabmyo-1.1.0}, we construct six conditions spanning channel count, sample rate, and bit depth. Every condition is derived from the same underlying recordings, so all comparisons are exactly paired at the subject level and no cross-dataset or cross-subject variability confounds the result. We then re-implement the resulting pipeline in C\texttt{++} on an ATmega328P microcontroller and measure what it costs to run.

This study addresses the following question: under a frozen edge deployment, which hardware constraint---spatial resolution, temporal rate, or signal fidelity---drives cross-session accuracy loss, and do these constraints interact with temporal drift? Throughout, we use \emph{compound} to denote a super-additive interaction: a cross-session accuracy loss under commodity constraints that exceeds the loss under clinical acquisition, rather than the hardware and temporal penalties simply summing.

The main contributions are as follows:

\begin{itemize}
    \item We decompose commodity hardware constraints into three independent axes and show that reducing sampling rate to \fsCommodity{}~Hz and bit depth to 10 bits changes cross-session degradation by less than 2 percentage points, while reducing channel count from 8 to 2 costs 11.10 percentage points of baseline accuracy and adds 9.09 percentage points of cross-session degradation.

    \item We show that the compounding of hardware and temporal effects is model-dependent, and that constrained hardware compresses the benefit of classifier selection from a 9.3-percentage-point spread on 8 channels to 2.3 percentage points on 2 channels.

    \item We use feature attribution to characterize the failure mechanism, finding that attribution profiles showed a cosine similarity of 0.978 across 28 days despite a 14.0-point accuracy drop, consistent with shifted inputs evaluated against frozen decision boundaries rather than with model reorganization.

    \item We report measured flash, SRAM, and per-stage latency for the full pipeline on an ATmega328P, establishing which architectures are physically deployable and where the latency ceiling lies.
\end{itemize}

The remainder of this paper is organized as follows. Section~2 reviews prior work on myoelectric pattern recognition, cross-session degradation, electrode shift, and hardware-constrained deployment, and identifies the gap this study addresses. Section~3 describes the dataset, degradation pipeline, evaluation protocol, and embedded profiling method. Section~4 reports results. Section~5 discusses their implications for hardware design, Section~6 states limitations, and Section~7 concludes.

% ============================================================
% RELATED WORK
% ============================================================
\section{Related Work}

This section reviews prior research relevant to cross-session myoelectric control on constrained hardware. We organize the discussion into four areas: the classical pattern recognition pipeline, cross-session and inter-day degradation, electrode shift and placement variability, and hardware-constrained embedded deployment. We then identify the gap this study addresses.

\subsection{Pattern Recognition for Myoelectric Control}
Myoelectric control relies on a highly standardized pattern recognition pipeline~\cite{ASGHARIOSKOEI2007275}. Hudgins et al.~\cite{204774} established the canonical time-domain feature set---mean absolute value, zero crossings, slope sign changes, and waveform length---which remains the foundation of modern classification owing to its computational efficiency and real-time viability~\cite{1206493}. Root mean square amplitude is a closely related and widely adopted addition to this family. Subsequent feature selection studies have demonstrated that most alternative time-domain features are highly redundant, justifying the continued use of a compact standard set~\cite{PHINYOMARK20127420} even as deep learning approaches explore learned representations~\cite{10.3389/fbioe.2020.00158}. Because the feature extraction and classification components of this pipeline are mature and widely reproduced, the barriers to continuous edge deployment are primarily physical and temporal rather than a question of which classifier to select.

\subsection{Cross-Session and Inter-Day Degradation}
Despite substantial advances in adaptive learning and domain generalization, cross-session and inter-day distribution shifts remain a central bottleneck limiting the long-term reliability of sEMG pattern-recognition systems in real-world applications. Inter-day signal non-stationarity causes progressive performance degradation as signal characteristics evolve over time through physiological factors and electrode shift~\cite{10.3389/fnbot.2018.00058}. Palermo et al., for example, observed an inter-session penalty on intact subjects in the NinaPro DB6 dataset~\cite{8009405}. The standard mitigation for this temporal drift relies on domain adaptation, adversarial self-calibration, or feature fusion to recover lost accuracy~\cite{doi:10.1142/S0219843620500255, 7302056, 10.3389/fnins.2017.00379, 8630679, 9207910}, though recent evidence questions the efficacy of adaptation across subjects and sessions~\cite{s21227500}.

Crucially, while these studies quantify temporal degradation carefully, these evaluations predominantly utilize multi-channel research-grade acquisition hardware. Although comparisons involving commodity and consumer-grade sEMG sensors have become increasingly common, many evaluations remain confined to within-session performance. Relatively few studies conduct extensive longitudinal assessments spanning multiple days or weeks under realistic deployment conditions. Consequently, cross-session robustness and long-term signal degradation remain major open challenges in real-world myoelectric control. Furthermore, the mitigations proposed in the longitudinal literature generally require dynamic model adaptation, retraining, or recalibration.

\subsection{Electrode Shift and Placement Variability}
A primary mechanism driving cross-session decay is physical shift of the electrode--skin interface after doffing and redonning the device. Hargrove et al. established that electrode displacements adversely affect classification accuracy, prompting the development of shift-aware training strategies~\cite{HARGROVE2008175}. Young et al. subsequently demonstrated that modifying electrode size and orientation can partially mitigate this sensitivity~\cite{5872014}, while others have applied transfer learning to counteract electrode displacement specifically~\cite{8673876}. While these studies correctly identify channel shift as a dominant failure mode, their proposed mitigations generally require redundant sensor arrays, subject-specific recalibration, or customized electrode geometries. This motivates our use of a rigid, landmark-anchored placement that reflects the generalized, non-optimized fit of commodity hardware.

\subsection{Hardware-Constrained and Embedded sEMG}
In parallel, a distinct body of literature has optimized sEMG classification for edge deployment on heavily constrained hardware. Li et al. evaluated the effect of sampling rate on classification accuracy to determine appropriate temporal constraints for prosthesis control~\cite{5626224}. Building on this, fully embedded systems have achieved real-time inference on multi-core biomedical SoCs~\cite{7974234}. More recent architectures have deployed deep learning, using temporal convolutional networks on multicore IoT processors~\cite{8930945} and in-sensor adaptive machine learning~\cite{moin2021wearable}, though these rely on processors orders of magnitude more capable than standard 8-bit microcontrollers such as the ATmega328P.

Pizzolato et al. directly compared six acquisition setups and found that commodity hardware performed comparably to research-grade equipment~\cite{Pizzolato2017-zd}. Our results are consistent with this finding; however, their evaluation used two Thalmic Myo armbands simultaneously (maintaining 16 channels and therefore introducing no spatial reduction), confounded amplifier and electrode differences with channel count, and evaluated performance within a single session. More recently, benchmark datasets such as SeNic have characterized non-ideal operating conditions on commodity hardware, including electrode shift, individual difference, inter-day variation, muscle fatigue, and arm posture~\cite{9771219}. That work establishes that these factors degrade recognition performance on a physical consumer device. However, the Myo armband used provides eight differential channels in a circumferential array, so it retains the spatial coverage of clinical systems. Evaluating a physical device also treats hardware constraints as a single aggregate variable, which precludes decomposing which specific constraint---spatial resolution, temporal rate, or signal fidelity---drives the failure.

\subsection{The Hardware--Time Gap}
The literatures addressing longitudinal sEMG degradation and embedded myoelectric control have largely evolved in isolation. Studies characterizing cross-session robustness predominantly rely on multi-channel research-grade acquisition~\cite{8009405, 7302056, pradhan2022grabmyo}. Conversely, research investigating commodity and microcontroller deployment either evaluates static models within a single session~\cite{7974234}, masks temporal drift through dynamic on-device retraining~\cite{moin2021wearable}, or evaluates the hardware as an aggregate physical device~\cite{9771219}. Consequently, it remains unclear whether specific hardware constraints interact with cross-session temporal drift under a frozen edge deployment, and which specific constraint---spatial resolution, temporal rate, or signal fidelity---drives any such interaction. This study addresses that gap by systematically degrading a clinical baseline toward a commodity edge configuration across a 28-day interval. Because all six experimental conditions derive from the identical underlying recordings, the resulting comparisons are exactly paired, isolating the source of performance loss without confounding cross-dataset or cross-subject variability.

% ============================================================
% METHODS
% ============================================================
\section{Methods}

\subsection{Dataset}
This study utilizes the complete \nSubjects{}-subject GRABMyo database, hosted on PhysioNet and distributed under a Creative Commons Attribution 4.0 International (CC-BY 4.0) license. Each able-bodied participant completed recording sessions across three separate days (Days \sessionDays).

Electromyographic signals were acquired using a forearm setup of 16 monopolar electrodes recorded via an EMGUSB2+ amplifier with a hardware gain of 500. The data were sampled at \fsClinical{} Hz and pre-processed by the dataset authors using a 10--500 Hz Butterworth bandpass filter and a 60 Hz notch filter. The sensor array consisted of two parallel circumferential rings spaced \interRingCm{} cm apart along the forearm axis. The proximal array (Ring 1) was longitudinally positioned at one-third the forearm length from the olecranon, with electrode 1 angularly aligned to the centerline of the elbow crease.

To evaluate a fundamental control scheme spanning the anterior and posterior compartments, we extracted three classes: wrist extension (Dataset Gesture 11, remapped to Class 0), hand close (Dataset Gesture 16, remapped to Class 1), and rest (Dataset Gesture 17, remapped to Class 2). Each gesture was executed across 7 trials per session, with each trial lasting 5 s. To eliminate physical transient onset and offset artifacts, the first and last 0.5 s of each trial were trimmed prior to filtering.

During initial data validation, we accounted for two discrepancies between the repository contents and the official documentation: inconsistent session folder naming conventions and the undocumented inclusion of the 17th gesture class.

\subsection{Hardware Degradation Pipeline} \label{sec:hardware_degradation}
To isolate the mechanical and signal-processing drivers of cross-session accuracy loss, the clinical baseline data were systematically degraded across three hardware axes: spatial resolution (channel count), temporal resolution (sample rate), and signal fidelity (bit depth). This produced six precisely controlled experimental conditions, summarized in Fig.~\ref{fig:pipeline}:

\begin{itemize}
    \item \textbf{Clinical (Baseline):} 8 bipolar channels, \fsClinical{} Hz sampling, unquantized floating-point precision.
    \item \textbf{Rate-only:} 8 channels, downsampled to \fsCommodity{} Hz, floating-point.
    \item \textbf{Depth-only:} 8 channels, \fsClinical{} Hz, 10-bit quantization.
    \item \textbf{Rate-depth:} 8 channels, \fsCommodity{} Hz, 10-bit quantization.
    \item \textbf{Spatial-only:} 2 channels (F1F9 and F5F13), \fsClinical{} Hz, floating-point.
    \item \textbf{Commodity:} 2 channels (F1F9 and F5F13), \fsCommodity{} Hz, 10-bit quantization.
\end{itemize}

\begin{figure}[htbp]
    \centering
    \includegraphics[width=\linewidth]{figures/fig_pipeline.pdf}
    \caption{Processing and evaluation pipeline. Records are loaded, converted
    to bipolar derivations, degraded along the hardware axes, filtered
    causally, and segmented into windowed features. The resulting feature
    cache supports both offline statistical evaluation and a separate
    re-implementation profiled on an ATmega328P microcontroller. All six
    conditions derive from the same underlying recordings.}
    \label{fig:pipeline}
\end{figure}

Signal degradation was applied in a strict sequential chain to mimic physical embedded processing constraints. Because all conditions were generated from the identical underlying recordings, every statistical comparison in this study is exactly paired at the subject level. Simulated quantization preceded the digital high-pass filter implemented in the target pipeline, reflecting the order in which these operations occur on the deployed device. The source recordings had already undergone the dataset's acquisition and preprocessing filters. When rate degradation was required, the signal was downsampled using a rational resampling ratio (125/256), which inherently applied an anti-aliasing FIR filter. When depth degradation was required, the signal was mapped to a symmetric bipolar 10-bit resolution based on a predefined maximum voltage threshold ($V_{MAX}$).

The $V_{MAX}$ thresholds were calibrated to the 99.9th percentile of the differenced, trimmed signal at the exact point in the processing chain where each respective quantizer operates. To prevent longitudinal test-set leakage, this calibration was restricted to the complete Session 1 dataset, explicitly excluding all cross-session test data. Importantly, the $V_{MAX}$ thresholds for the 8-channel depth ablations (\vmaxEightCh{} mV for \texttt{depth\_only} and \vmaxEightChResampled{} mV for \texttt{rate\_depth}) were pooled across the entire 8-channel array. This effectively quantized the primary channels roughly 1.9 times more coarsely than the 2-channel commodity condition ($V_{MAX} = $ \vmaxTwoCh{} mV). This asymmetry was preserved intentionally to act as a conservative handicap, ensuring any robustness demonstrated under the depth limitation is mathematically bounded.

\subsection{Channel Selection and Bipolar Derivation}
To mimic the physical hardware of a 2-channel commodity edge sensor, monopolar recordings were converted to bipolar derivations. Adjacent electrodes between Ring 1 and Ring 2 were subtracted (e.g., F1 minus F9).

The 2-channel commodity condition isolated a single flexor-extensor pair spaced 180 degrees apart (four 45-degree steps on the circular array) to capture the anterior and posterior compartments of the forearm. Because hand close is a flexor-dominant gesture and wrist extension is extensor-dominant, this mapping provides the physiological basis for baseline separability. This landmark-anchored placement deliberately avoids per-subject spatial optimization, accurately reflecting the rigid, generalized fit of consumer armbands. To ensure strict spatial parity, the 2 commodity channels (F1F9 and F5F13) operate as a direct subset of the 8 clinical channels (Derivations 0 and 4).

In a representative trial, prior to subtraction, adjacent monopolar channels (e.g., F1 and F9) exhibited a high Pearson correlation ($r =$ \rawCorr), indicating heavy common-mode content and justifying spatial high-pass rejection via differencing. Following subtraction, the resulting bipolar channels (F1F9 and F5F13) were substantially decorrelated ($r =$ \bipolarCorr), suggesting that they carry complementary spatial information.

\subsection{Filtering and Feature Extraction}
A causal 4th-order Butterworth high-pass filter (\hpCutoff{} Hz) was applied to remove baseline wander and motion artifacts. Zero-phase filtering techniques were strictly avoided to maintain causality, ensuring the pipeline remains viable for real-time edge deployment. No additional low-pass filter was applied beyond the dataset's native 500 Hz acquisition limit. Because the source recordings were already band-limited to 500 Hz, the \fsCommodity{}~Hz condition preserves the nominal source passband at its Nyquist limit. The rate ablation therefore primarily evaluates reduced temporal sampling density rather than a broad bandwidth reduction, with practical attenuation near the upper band edge set by the anti-aliasing response of the polyphase resampler.

Following the filter chain, an additional \transientMs{} ms trim was discarded to eliminate transient IIR filter startup artifacts. This trim was floored to sample indices to maintain exact window-count symmetry across all resampling arms. Signals were then segmented into \windowMs{} ms windows with a \stepMs{} ms (50\%) step overlap, yielding exactly \nWindows{} windows per trial (totaling \nRows{} windows per condition).

Four standard time-domain features were extracted per channel: Root Mean Square (RMS), Mean Absolute Value (MAV), Waveform Length (WL), and Zero Crossings (ZC). RMS and MAV are inherently rate-independent. To ensure feature parity across the temporal degradation conditions, WL was divided by the window sample count, and ZC was divided by the window duration (yielding crossings per second). A \zcEpsilon{} $\mu$V deadzone was applied to the ZC calculation to prevent hardware noise inflation.

\subsection{Evaluation Protocol}
The classification task was structured as a subject-dependent, longitudinal evaluation. To ensure strict training parity, models for all six degradation conditions were trained exclusively on Session 1 (Day 1), Trials 1--5. Baseline performance was established using Session 1, Trials 6--7. Degradation curves were tracked through Session 2 (Day 8), and maximum cross-session robustness was evaluated against Session 3 (Day 29), Trials 6--7.

To prevent data leakage across the 50\% window overlap, the independent unit of observation was strictly maintained at the trial level. Feature standardization was applied using a scaler fit exclusively on the training partition. The classification pipeline evaluated five models: Linear Discriminant Analysis (LDA), LinearSVC, K-Nearest Neighbors ($k=5$), Random Forest (50 estimators), and XGBoost (50 estimators). All stochastic algorithms (Random Forest, XGBoost, LinearSVC) and bootstrap resampling procedures utilized a fixed random seed (42) to ensure strict reproducibility. Data splits, LDA, and KNN relied entirely on deterministic operations.

\subsection{Deployability Profiling Method}
SRAM and Flash memory footprints were profiled in C++ to determine deployability constraints for a standard ATmega328P microcontroller (16 MHz, 2 KB SRAM, 32 KB Flash). A baseline hardware profiling sketch established the memory budget available for model parameters. Footprints were calculated for the commodity target (2 channels, 8 features, 3 classes) using exact float32 parameter counts representing the physical device payload, rather than serialized file sizes. All deployment structures include the fitted standard scaler (16 floats, 64 bytes). Footprints for fixed-architecture models (LDA, LinearSVC, KNN) were profiled using a representative subject, while tree-based ensembles (Random Forest, XGBoost) report the minimum and maximum footprint across all \nSubjects{} subjects due to their data-dependent structural variance.

Model payloads were defined per architecture: weight and bias vectors for LDA and LinearSVC, the complete training feature matrix and its corresponding labels for KNN (i.e., 555 windows $\times$ 8 features $\times$ 4 bytes, plus 555 label bytes), and explicit node array structures for tree-based models. Tree layouts assumed 9 bytes for internal nodes (1-byte feature index, 4-byte float threshold, 4-byte child pointers), 1 byte for Random Forest leaves (class label), and 4 bytes for XGBoost leaves (float weight).

To evaluate SRAM constraints, three mutually exclusive buffer architectures were analyzed: a naive float32 window (1,600 bytes), an \texttt{int16} window (800 bytes), and a streaming accumulator architecture measured at \accumBytesMeasured{} bytes on the target. The streaming design accumulates features over \stepMs{} ms sub-blocks without storing the raw window, and is the architecture carried through to the verified implementation in Section~\ref{sec:embedded_verification}.

Microsecond-resolution execution timing was captured to measure ADC, filtering, and inference latencies against a 1 kHz sampling deadline. Per-stage latencies were measured using a loop-and-divide methodology, with single-shot \texttt{micros()} call overhead and empty-loop iteration overhead strictly subtracted from all reported measurements. To prevent the compiler's optimizer from silently deleting unconsumed arithmetic during benchmarking, a \texttt{volatile} sink was utilized. The worst-case execution time (WCET) was captured using a separate methodology: timing 1,000 individual window-boundary ticks and recording the maximum duration. These latencies were obtained from a dedicated profiling sketch that reproduces the computational workload of the deployed pipeline, using the same filter order, feature set, and window structure, but implemented with direct-form biquad sections and a single running accumulator. It is therefore distinct from the verified firmware described in Section~\ref{sec:embedded_verification}, which uses transposed direct form II and rotating self-contained sub-blocks. The reported latencies are representative of the per-stage computational cost rather than measurements taken on the verified implementation.

\subsection{Embedded Implementation Verification} \label{sec:embedded_verification}
To establish that the deployed implementation reproduces the offline analysis rather than merely resembling it, we performed a three-level numerical comparison. Level A is the analysis pipeline described above, using SciPy and scikit-learn in double precision. Level B is an explicit single-precision reference implemented in Python that mirrors the operation order of the embedded code: transposed direct form II filtering applied sample by sample, multiplication by a precomputed reciprocal in place of division during standardization, and ordered accumulation of the discriminant dot product. Level C is the physical ATmega328P. The A-to-B comparison isolates the numerical consequence of these deployment choices; the B-to-C comparison isolates implementation correctness.

Because the streaming architecture processes each sample once and cannot revisit earlier data, features are accumulated over \stepMs{} ms sub-blocks that are combined pairwise into \windowMs{} ms windows. Each sub-block retains only statistics computed from its own samples. The single sample-to-sample difference spanning the boundary between two adjacent sub-blocks belongs to neither and is added when the pair is joined; retaining it inside either block would double count it, because each sub-block contributes to two consecutive overlapping windows. This reformulation was verified against the direct windowed implementation on 3,500 synthetic cases constructed to exercise the zero-crossing deadzone and sign conventions at sub-block boundaries, and on real trials, with agreement to machine precision.

Quantized 10-bit codes were streamed to the microcontroller over a framed serial protocol with per-chunk acknowledgement, so transfer rate could not perturb the arithmetic. A CRC-16 computed over the received bytes confirmed that codes arrived intact, distinguishing transport faults from implementation error. Filter and accumulator state were reset at each trial boundary, matching the reference. The device returned raw and standardized features, discriminant scores, and predicted classes for every window; a separate diagnostic mode additionally returned every filtered sample for one trial.

Acceptance criteria were fixed before any hardware result was observed: exact agreement for ADC codes, zero-crossing counts, window counts, and predictions; filter output within 0.5 ADC least significant bits; standardized features within $10^{-3}$ training standard deviations; and discriminant scores within a relative tolerance of $10^{-4}$ and an absolute tolerance of $10^{-5}$. Because a per-score error of magnitude $E$ can lower the winning score and raise the runner-up simultaneously, the observed margin shrinks by $2E$; we therefore adopted $2E < M$, for minimum decision margin $M$, as a sufficient worst-case condition for preserving the predicted class. It is not a necessary condition, and predictions may agree when it fails. A comparison of the continuous raw features was added after the initial hardware results were observed; it is reported as a post-hoc diagnostic and did not form part of the acceptance criteria.

\subsection{Statistical Analysis}
To separate exploratory hypothesis generation from rigorous verification, the statistical analysis plan was formally locked after an initial observation of the primary main effects. The evaluation comprises three distinct analyses.

Four planned paired Wilcoxon contrasts characterized the LDA hardware-by-session cell means: clinical versus commodity on Day 1, clinical versus commodity on Day 29, Day 1 versus Day 29 under clinical acquisition, and Day 1 versus Day 29 under commodity constraints. A Bonferroni-adjusted threshold of $\alpha = 0.0125$ was applied.

The primary analysis evaluates the model dependence of the cross-session drop (the 2x2 interaction). The performance of the clinical baseline vs. commodity hardware across Sessions 1 and 3 was evaluated for all five model architectures. Significance was tested using a two-sided paired Wilcoxon signed-rank test. Confidence intervals (95\%) for the interaction were computed using a percentile bootstrap over 10,000 resamples. To correct for multiple comparisons across the five models, a Bonferroni threshold of $\alpha = \bonferroniAlpha$ was applied.

A secondary descriptive analysis evaluated this cross-session drop normalized against available performance headroom (Day 1 accuracy minus a 33.3\% theoretical chance floor). To prevent division by near-zero, a union exclusion mask was applied: any subject scoring $\le$ 40\% on Day 1 under any model in either condition was dropped for this specific endpoint ($n=41$). The 40\% threshold was chosen to place the denominator at least 6.7 percentage points above the 33.3\% chance floor; it is a pragmatic bound rather than a principled one, and it applies only to this secondary endpoint. Unlike the primary metric, this descriptive endpoint was not subject to multiplicity correction.

The third analysis, the hardware ablation, isolates the precise source of the performance penalty. The independent endpoints are the difference in cross-session degradation (Day 1 minus Day 29 accuracy) and the difference in Day 1 baseline accuracy between the clinical baseline and each of the five degradation variants. To bound the exact effect sizes without discarding magnitude, 95\% confidence intervals for the paired per-subject differences were calculated using the Bias-Corrected and accelerated (BCa) bootstrap method. No multiplicity correction was applied to the ablation intervals because this analysis is estimation-focused. The intervals are interpreted individually and descriptively and should not be treated as simultaneous family-wise bounds.

\subsection{Feature Attribution Method}
To evaluate the spatial redundancy hypothesis mechanically, Random Forest feature attribution was analyzed across all six degradation conditions using TreeExplainer (SHAP)~\cite{lundberg2017unifiedapproachinterpretingmodel}. The analysis isolated the eight features originating from the physical locations shared between devices (F1F9 and F5F13) to measure two metrics specified in the locked analysis plan between Day 1 and Day 29: Attribution Migration (the fraction of total model decision weight placed on the shared channels) and Profile Shift (the cosine distance of the normalized 8-feature attribution profile).

To isolate true longitudinal shift from dimensionality noise, cross-session profile shifts were normalized against an empirical within-day baseline. The baseline averaged within-day noise (Day 1 Trial 6 vs. Trial 7, and Day 29 Trial 6 vs. Trial 7), while the longitudinal shift averaged cross-day matched trials (T6 vs. T6, and T7 vs. T7). Because all four profiles were computed from single trials (37 windows each), window counts were perfectly matched. The four 8-channel conditions served as a falsification control: the analysis plan stated that if channel count drove the effect, the 8-channel models would show significantly less profile shift than the 2-channel models, and attribution migration would decrease as the shared features became less discriminative on Day 29 against the frozen training-time splits.

% ============================================================
% RESULTS
% ============================================================
\section{Results}

\subsection{Baseline Performance}
The clinical within-session evaluation achieves 96.2\% accuracy under LDA and 97.2\% under Random Forest. This falls within the range reported for comparable time-domain classification on multi-channel forearm sEMG; the putEMG benchmark, for instance, reports approximately 90\% for LDA and SVM on an 8-gesture task~\cite{s19163548}. Our task evaluates three classes and eight bipolar derivations using the open-access GRABMyo dataset. Because this configuration is narrower than the dataset's full gesture set and electrode array, this baseline describes our observed system performance rather than reproducing published metrics.

\subsection{Hardware and Session Effects}
\begin{figure*}[t]
    \centering
    \includegraphics[width=\textwidth]{figures/confusion_matrices.pdf}
    \caption{Cross-session confusion matrices for the Clinical (8-channel) and Commodity (2-channel) conditions on Day 1 and Day 29. Error is concentrated in Extension$\leftrightarrow$Close confusion across all conditions. Hardware degradation scales this confusion approximately threefold without altering its fundamental structure. Rest recall is nearly unaffected by session change under clinical hardware (97.4\% to 96.4\%) but falls from 96.3\% to 82.8\% under commodity constraints, with \cmRestMisclass{} of Rest windows classified as gestures.}
    \label{fig:confusion}
\end{figure*}

The core performance limits are defined by the primary experimental axes rather than algorithmic selection. Figure~\ref{fig:confusion} illustrates the two-by-two interaction of hardware constraints and cross-session temporal drift. Under LDA classification, the system achieves 96.2\% for clinical within-session, 86.1\% for commodity within-session, 82.5\% for clinical cross-session, and 71.4\% for commodity cross-session.

All four planned pairwise contrasts were statistically significant after Bonferroni correction for four tests ($\alpha = 0.0125$) via paired Wilcoxon tests. Within a single session, degrading the system from clinical to commodity hardware yields a significant accuracy penalty ($p = 3.6 \times 10^{-8}$). This hardware penalty persists cross-session ($p = 1.7 \times 10^{-3}$). Similarly, the introduction of a 28-day inter-session gap significantly impacts performance under both the clinical baseline ($p = 2.7 \times 10^{-4}$) and the fully degraded commodity constraints ($p = 1.1 \times 10^{-4}$).

\subsection{Model Dependence}
Algorithmic architecture dictates whether hardware constraints act additively or compound cross-session performance loss. To evaluate this interaction, we compared the clinical and commodity degradation slopes across the five tested models. Table \ref{tab:model_dependence} reports the absolute accuracy drop, the headroom-normalized loss, paired Wilcoxon tests, and 95\% percentile bootstrap confidence intervals.

\begin{table}[htbp]
    \centering
    \small
    \setlength{\tabcolsep}{4pt}
    \caption{Cross-session degradation (Day 1 to Day 29) by model architecture. Clinical and commodity columns give the raw accuracy drop in percentage points; the difference column gives the paired interaction with a 95\% percentile bootstrap interval over 10,000 resamples. $p$ values are paired Wilcoxon tests. $p_{\text{norm}}$ (italicized) is the headroom-normalized secondary endpoint ($n=41$ after the union exclusion mask), which is descriptive and uncorrected for multiplicity.}
    \label{tab:model_dependence}
    \begin{tabular}{lcccrr}
        \toprule
        \textbf{Model} & \textbf{Clin.} & \textbf{Comm.} &
        \textbf{Diff. [95\% CI]} & \textbf{$p$} & \textit{$p_{\text{norm}}$} \\
        \midrule
        LDA           & 13.7 & 14.7 & $+1.0$ [$-5.3$, $+7.2$]  & 0.5869 & \textit{0.4801} \\
        KNN           & 13.7 & 11.6 & $-2.1$ [$-9.5$, $+5.1$]  & 0.7082 & \textit{0.8275} \\
        LinearSVC     & 10.5 & 16.8 & $+6.3$ [$-0.6$, $+13.0$] & 0.0720 & \textit{0.0265} \\
        Random Forest & 5.4  & 14.0 & $+8.6$ [$+2.8$, $+14.5$] & 0.0091 & \textit{0.0049} \\
        XGBoost       & 9.0  & 14.7 & $+5.7$ [$+0.3$, $+11.5$] & 0.0848 & \textit{0.0285} \\
        \bottomrule
    \end{tabular}
\end{table}

The compounding effect of hardware degradation is highly model-dependent. We did not detect a compounding interaction for the linear discriminant or nearest-neighbor classifiers. However, their confidence intervals are wide: for LDA, the interaction is bounded between $-5.3$ and $+7.2$ points, so effects as large as 7 points remain consistent with the data. Rather than proving additivity, these tests are underpowered to detect compounding. In contrast, the Random Forest exhibits a significant interaction and is the only model to survive a conservative Bonferroni correction ($\alpha = 0.01$).

Notably, the XGBoost absolute difference confidence interval strictly excludes zero ($[+0.3, +11.5]$), yet the paired Wilcoxon test fails to reject the null ($p=0.0848$). This divergence indicates right-skewed paired differences: a minority of subjects with large positive differences pulls the bootstrap mean upward, while the signed-rank test provides only weak evidence of a systematic location shift. LinearSVC exhibits a milder iteration of this skew.

Consequently, model choice yields diminishing returns under hardware constraints. A post-hoc observation of Day 29 accuracies illustrates this compression: algorithmic selection accounts for a 9.3 percentage-point spread in the clinical condition, but only a 2.3 percentage-point spread in the commodity condition. Structurally, the interaction manifests because the cross-session drop on commodity hardware remains relatively uniform across all models (11.6 to 16.8 points), whereas the clinical drop varies widely (5.4 to 13.7 points). The Random Forest does not degrade faster on commodity hardware so much as it degrades unusually little on clinical hardware. One candidate explanation, examined further in Section~\ref{sec:model_dependence_discussion}, is training-time spatial redundancy. With eight channels available, the ensemble distributes splits across them, meaning drift in any single channel invalidates a smaller fraction of the frozen trees. Under the 2-channel spatial ablation of the commodity condition, this baked-in redundancy is eliminated.

\subsection{Degradation Axis Attribution}
To isolate the precise hardware drivers of cross-session accuracy loss, the clinical baseline was systematically degraded across three hardware axes: temporal (rate), signal fidelity (depth), and spatial (channel count). Table \ref{tab:ablation} and Fig.~\ref{fig:ablation} present the isolated effect of each degradation axis on both Day 1 baseline accuracy and cross-session degradation.

\begin{table}[htbp]
    \centering
    \caption{Degradation axis attribution. Paired per-subject differences from the clinical baseline (Random Forest, $n=\nSubjects{}$), with 95\% BCa bootstrap confidence intervals. Negative Day 1 values indicate lower baseline accuracy than clinical; positive drop values indicate steeper cross-session degradation than clinical.}
    \label{tab:ablation}
    \begin{tabular}{lcc}
        \toprule
        \textbf{Condition} & \textbf{Day 1 diff (pts)} & \textbf{Drop diff (pts)} \\
        \midrule
        Rate-only    & \rateDayOneDiff{} \rateDayOneCI{}         & \rateDropDiff{} \rateDropCI{} \\
        Depth-only   & \depthDayOneDiff{} \depthDayOneCI{}       & \depthDropDiff{} \depthDropCI{} \\
        Rate-depth   & \rateDepthDayOneDiff{} \rateDepthDayOneCI{} & \rateDepthDropDiff{} \rateDepthDropCI{} \\
        \midrule
        \textbf{Spatial-only} & \textbf{\spatialDayOneDiff{} \spatialDayOneCI{}} & \textbf{\spatialDropDiff{} \spatialDropCI{}} \\
        Commodity    & \commodityDayOneDiff{} \commodityDayOneCI{} & \commodityDropDiff{} \commodityDropCI{} \\
        \bottomrule
    \end{tabular}
\end{table}

\begin{figure}[htbp]
    \centering
    \includegraphics[width=\linewidth]{figures/fig_ablation.pdf}
    \caption{Degradation axis attribution. Paired per-subject differences from the clinical baseline with 95\% BCa bootstrap confidence intervals (Random Forest, $n=\nSubjects{}$). Sample rate and bit depth reductions (grey) are bounded near zero on both endpoints, while channel reduction (red) shifts both baseline accuracy and cross-session degradation substantially. Note that the sign conventions are opposite between panels.}
    \label{fig:ablation}
\end{figure}

The ablation reveals that temporal and signal fidelity constraints contribute negligibly over the ranges tested. The 95\% confidence intervals for the \texttt{rate-only}, \texttt{depth-only}, and \texttt{rate-depth} conditions are tightly bounded around zero. Specifically, reducing the sampling rate to \fsCommodity{} Hz and bit depth to 10-bit quantization alters cross-session degradation by less than 2 percentage points and Day 1 accuracy by less than 0.75 points. The \texttt{depth-only} Day 1 interval reaches zero at the reported precision, while its point estimate of 
$-0.09$ points is practically negligible. Furthermore, as outlined in Section~\ref{sec:hardware_degradation}, this depth ablation was conservatively estimated using a 1.9$\times$ coarser pooled $V_{MAX}$ threshold, reinforcing the system's overall robustness to quantization.

In contrast, spatial resolution dictates performance. The \texttt{spatial-only} reduction from eight to two channels drops Day 1 baseline accuracy by 11.10 points and accelerates cross-session degradation by an additional 9.09 points compared to the clinical baseline. Crucially, the direct paired difference between the \texttt{commodity} and \texttt{spatial-only} conditions is $+0.87$ points $[-0.67, +2.42]$ for Day 1 baseline accuracy, and $-0.52$ points $[-3.06, +1.26]$ for cross-session degradation. This indicates that adding rate and depth degradation to an already channel-limited system extracts no further measurable penalty, establishing channel count as the dominant hardware constraint within the tested ranges.

Most of the observed cross-session decline was already present by Day 8 in both conditions, though the magnitudes differ substantially. Under Random Forest, clinical accuracy was 97.2\%, 91.5\%, and 91.8\% on Days 1, 8, and 29, while commodity accuracy was 86.9\%, 72.1\%, and 73.0\%. With only three measurement sessions, these results do not establish when within the first week the decline occurred, but they motivate denser evaluation during the early deployment period.

\subsection{Attribution Analysis}
To determine whether hardware-induced accuracy loss stems from the model fundamentally changing its decision logic or from degraded signals shifting against static thresholds, we analyzed feature attribution profiles. While both hypotheses specified in the locked analysis plan yielded statistically null results, the profile shift direction was consistent with predictions.

Feature-importance migration showed no statistically detectable change for the clinical condition, shifting from 15.0\% to 14.9\% ($p = 0.92$). (Migration for the 2-channel conditions is pinned at 100\% by construction). For the spatial attribution shift, all four 8-channel conditions (0.0041, 0.0042, 0.0048, 0.0067) exhibited lower profile shift distances than both 2-channel conditions (0.0086, 0.0131). This aligns with the directional prediction stated in the analysis plan, though the \texttt{depth-only} condition (0.0067) sat near the boundary, and none of the eight comparisons survived a Bonferroni correction at $\alpha = 0.0063$.

Notably, the eight features derived from the two shared channels represent 25\% of the clinical feature set but receive only 15\% of the model's total attribution. This indicates that the landmark-chosen forearm channels are below-average in informational value when alternatives are available, aligning with the $V_{MAX}$ calibration, in which F1--F9 exhibited the lowest 99.9th-percentile amplitude of the eight bipolar derivations (\fOneFNineVmax{} mV, versus \fSixFFourteenVmax{} mV for the largest, F6--F14).

Commodity cross-session attribution profiles showed a cosine similarity of 0.978 despite a 14.0-point drop in classification accuracy. This is consistent with shifted inputs being evaluated against fixed decision boundaries, rather than with the model altering its internal feature weighting.

\subsection{Deployability}

\textbf{Dynamic Memory (SRAM)} \\
A naive \windowMs{} ms window of 32-bit floats at a \fsCommodity{} Hz sampling rate requires a 1,600-byte data buffer, which alone consumes 78\% of the 2,048-byte capacity and leaves insufficient room for the Arduino runtime and serial buffers. We avoided this with the streaming accumulator described in Section~\ref{sec:embedded_verification}, which stores no raw window at all. Measured on the target with \texttt{sizeof}, each sub-block occupies \subBlockBytes{} bytes; two rotating sub-blocks per channel across two channels total \accumBytesMeasured{} bytes. Filter state adds \filterStateBytes{} bytes, giving \dspStateBytes{} bytes of persistent signal-processing state. The complete verification firmware, including the classifier parameters, the streaming feature extractor, and the serial harness, compiles to \verifySRAM{} bytes of static SRAM (23\% of capacity) and \verifyFlash{} bytes of flash (32\%).

\textbf{Program Storage (Flash)} \\
The base pipeline, including the filter, accumulators, and profiling test harness, consumes \measuredFlash{} bytes of Flash. This establishes a conservative \flashBudget{}-byte budget for model parameters stored in PROGMEM. Profiling reveals that edge deployability is heavily model- and subject-dependent (Table \ref{tab:memory}). While parametric models fit easily, tree-based models scale with subject variance. The Random Forest architecture fits for the subject yielding the smallest model but exceeds the Flash limit at 117\% capacity for the subject yielding the largest model.

\begin{table}[htbp]
    \centering
    \caption{Memory footprint and deployability verdict per model architecture. Total flash includes the \measuredFlash{}-byte base pipeline. \emph{min} and \emph{max} denote the subjects yielding the smallest and largest models.}
    \label{tab:memory}
    \begin{tabular}{lrrl}
        \toprule
        \textbf{Model} & \textbf{Payload (B)} & \textbf{Total flash (B)} & \textbf{Verdict} \\
        \midrule
        LDA               & 172    & 8,516 (26\%)   & Fits \\
        LinearSVC         & 172    & 8,516 (26\%)   & Fits \\
        KNN ($k=5$)       & 18,379 & 26,723 (83\%)  & Fits \\
        RF (min)          & 1,334  & 9,678 (30\%)   & Fits \\
        XGBoost (min)     & 2,172  & 10,516 (33\%)  & Fits \\
        XGBoost (max)     & 20,957 & 29,301 (91\%)  & Fits \\
        \textbf{RF (max)} & \textbf{29,334} & \textbf{37,678 (117\%)} & \textbf{Does not fit} \\
        \bottomrule
    \end{tabular}
\end{table}

\textbf{Compute Timing and Architecture} \\
Maintaining a 1 kHz sampling schedule imposes a 1,000 $\mu$s deadline between samples. Table \ref{tab:latency} breaks down the pipeline execution cost. The per-sample ADC, filtering, and accumulation workload requires approximately 654 $\mu$s at the default ADC prescaler (PS=128), consuming 65.4\% of the sampling interval and leaving 34.6\% headroom before window-boundary processing.

However, the worst-case execution time (WCET) occurs at the window boundary, where finalization and LDA inference are added to the workload. The measured WCET at PS=128 is \wcetTick{} $\mu$s. Because this boundary tick exceeds the 1,000 $\mu$s limit, a sequential loop architecture will fail and drop samples every 100 ms.

\begin{table}[htbp]
    \centering
    \caption{Execution latency and CPU duty cycle breakdown, measured on the
    ATmega328P at the default ADC prescaler (PS=128). Latency in $\mu$s per
    invocation; rate in Hz; duty cycle as a percentage of available CPU time.}
    \label{tab:latency}
    \begin{tabular}{lrrr}
        \toprule
        \textbf{Stage} & \textbf{Latency} & \textbf{Rate} & \textbf{Duty} \\
        \midrule
        ADC read              & 223.79 & 1,000 & 22.4 \\
        Filter + accumulation & 430.68 & 1,000 & 43.1 \\
        Feature finalization  & 234.70 & 10    & 0.2 \\
        LDA inference         & 608.58 & 10    & 0.6 \\
        \midrule
        \textbf{System average} & & & \textbf{66.3} \\
        \bottomrule
    \end{tabular}
\end{table}

The inference calculation itself does not need to complete within 1 ms; it only needs to finish before the next 100 ms window closes. Therefore, successful edge deployment requires an interrupt-driven architecture, decoupling the strict 1 kHz ADC sampling routine from the asynchronous main loop evaluating the model. We also tested accelerating the ADC (PS=16) to gain an additional 192 $\mu$s of headroom per sample. While this trades 10-bit accuracy for speed (operating at 1 MHz, effectively reducing bit resolution), our offline \texttt{depth\_only} ablation showed that 10-bit quantization costs under 0.3 points of accuracy. This suggests the pipeline is relatively insensitive to quantization, making the prescaler tradeoff plausibly cheap, though the exact effective bit resolution was not characterized.

\textbf{System Latency} \\
Total end-to-end latency is 200 ms of signal collection plus $\approx$1.6 ms of computation, yielding approximately \latencyTotal{} ms. This exceeds the 100--125 ms optimal controller delay reported by Farrell and Weir~\cite{4126535}. However, this excess is attributable almost entirely to the analysis window, with computation accounting for less than 1\% of the total delay. Reducing latency therefore requires shortening the data collection window rather than accelerating inference, representing a strict tradeoff against feature stability rather than processor capability.

\textbf{Implementation Verification} \\
The three-level comparison was performed on the complete held-out partition for one subject: three sessions, three gestures, and two trials per condition, yielding \nVerifyWindows{} windows across \nVerifyTrials{} trials and spanning the full 28-day interval. Levels A and B agreed on all \nVerifyWindows{} predictions with no zero-crossing mismatches, a maximum standardized-feature difference of $\abStdErr$ standard deviations, and a maximum filter difference of $\abFiltErr$ least significant bits. The maximum per-score difference was $\abScoreErr$ against a minimum decision margin of $\abMargin$, satisfying $2E < M$.

Level C reproduced Level B on every predetermined criterion. Across all \nVerifyTrials{} trials, every ADC-code checksum matched, all \nVerifyWindows{} windows were produced, and zero-crossing counts, standardized features, discriminant scores, and predicted classes showed zero deviation. A separate sample-level comparison on one \nFilterSamples{}-sample trial recorded zero deviation in filter output. A post-hoc comparison of the continuous raw features, added after the initial results were observed and not among the acceptance criteria, likewise showed zero deviation. Single-precision arithmetic on the target device therefore introduced no measurable deviation from the reference at any stage examined.

% ============================================================
% DISCUSSION
% ============================================================
\section{Discussion}

This study systematically degraded a clinical sEMG baseline to isolate the precise hardware constraints driving cross-session accuracy loss. By evaluating identical underlying source recordings across a 28-day interval, we established an optimization hierarchy for hardware design, characterized the stability of feature-attribution patterns across sessions, and defined the physical deployability limits of microcontroller-based pattern recognition.

\subsection{Hardware Design and the Spatial Bottleneck}
When engineering a heavily constrained edge myoelectric classifier, the hardware ablation isolates spatial resolution as the dominant physical bottleneck. Degrading the temporal rate and signal fidelity to commodity levels (downsampling to \fsCommodity{} Hz and 10-bit quantization) altered cross-session degradation by less than 2 percentage points and Day-1 accuracy by less than 0.75 points, both with tightly bounded confidence intervals. Conversely, reducing the spatial coverage from eight channels to two landmark-anchored derivations cost 11.10 percentage points of baseline accuracy and added 9.09 percentage points to the Day-1-to-Day-29 cross-session drop.

Crucially, the axes do not combine additively. The direct paired difference between the fully degraded commodity condition and the spatial-only condition was bounded near zero ($+0.87$ points $[-0.67, +2.42]$ on Day 1; $-0.52$ points $[-3.06, +1.26]$ for the cross-session drop). Adding temporal and depth degradation to an already channel-limited system exacts no further penalty. For hardware designers, this suggests a clear allocation principle: within the ranges tested, resources are better spent on electrode count and placement stability than on high-resolution analog-to-digital converters or elevated sampling frequencies.

\subsection{Training-Time Redundancy and Model Dependence}\label{sec:model_dependence_discussion}
The compounding effect of these hardware limits on temporal drift is model-dependent. We did not detect a compounding interaction for simple linear and distance-based architectures (LDA, KNN). However, their confidence intervals were 12 to 15 points wide, meaning interaction effects up to 7 points remain consistent with the data.

In contrast, the Random Forest architecture exhibited a significant compounding interaction ($p = 0.0091$) that survived Bonferroni correction. One explanation for this interaction is training-time spatial redundancy. Tree ensembles distribute decision splits across available channels; when eight channels are present, temporal drift in any single channel invalidates a smaller fraction of the frozen trees. Under the two-channel commodity constraint, this spatial redundancy is eliminated.

An alternative explanation is that Random Forest simply generalizes better in general, independent of spatial redundancy: it attains the highest clinical Day-1 accuracy (\rfClinDayOne{}) and therefore has the most headroom to lose. The headroom-normalized secondary endpoint partially addresses this, since it divides each drop by the accuracy available above chance; Random Forest remains the strongest interaction under that normalization ($p = \rfInteractionNormP$). We cannot fully separate capacity from redundancy with the present design, and a channel-count sweep holding model architecture fixed would be required to do so. Consequently, heavily constrained hardware removes the ability to algorithmically engineer around the physical problem. A post-hoc observation of Day 29 accuracies illustrates this compression: algorithmic selection buys a 9.3 percentage-point spread in robustness on clinical hardware (82.5\%--91.8\%), but only a 2.3 percentage-point spread on commodity hardware (70.7\%--73.0\%).

\subsection{Timescale of Cross-Session Loss}
Most of the cross-session decline was already present at Day 8, and both conditions were approximately stable between Day 8 and Day 29 (clinical \rfClinDayEight{} to \rfClinDayTwentyNine{}; commodity \rfCommDayEight{} to \rfCommDayTwentyNine{}). The observed means do not show continued monotonic decline between Days 8 and 29, which is consistent with a largely one-time perturbation introduced when the electrode array is removed and reapplied, though drift that saturates early would produce a similar pattern. Three sessions cannot distinguish a doffing-and-redonning event from drift that saturates within the first week, and the sessions differ in both elapsed time and number of reapplications. Denser sampling during the first week, or repeated reapplication within a single day, would be required to separate these explanations. Practically, the observation suggests that recalibration schedules for spatially constrained devices should be evaluated on a scale of days rather than weeks.

\subsection{Attribution Stability Across Sessions}
The feature attribution analysis is consistent with a physical rather than algorithmic explanation for the observed accuracy loss. Despite a 14.0-point drop in classification accuracy, the Random Forest attribution profiles for the commodity condition showed a cosine similarity of 0.978 between Day 1 and Day 29. The longitudinal migration of attribution weight also showed no statistically detectable change ($p = 0.92$).

The model therefore relies on the same features in approximately the same proportions while producing more errors. This pattern is consistent with shifted feature values being evaluated against fixed training-time decision boundaries, although attribution stability alone does not establish the physical cause of the shift.

\subsection{Deployability and Latency Ceilings}
Profiling the models for an ATmega328P microcontroller underscores the deployability advantage of the compressed model spread. While the Random Forest provides superior clinical robustness, for the worst-case subject its 29,334-byte parameter payload combined with the \measuredFlash{}-byte base pipeline totals 37,678 bytes, exceeding the 32,256-byte usable flash limit at 117\% capacity. Conversely, the LDA model consumes only 172 bytes.

Because the commodity hardware constraint compresses the algorithmic variance, this memory tradeoff becomes highly favorable. Deploying LDA instead of a Random Forest on eight channels would cost 9.3 points of accuracy; on the two-channel commodity configuration, the same choice costs only 1.6 points.

However, while inference footprints are manageable, system latency remains a strict barrier. The end-to-end processing latency was measured at \latencyTotal{} ms, exceeding the 100--125 ms optimal controller delay reported by Farrell and Weir~\cite{4126535}. Because computation accounts for only 1.6 ms of this total, the delay is fundamentally a windowing problem. Overcoming this latency ceiling requires truncating the \windowMs{} ms data collection window, posing a direct tradeoff between real-time responsiveness and feature stability.

\subsection{Practical Implications and Prior Work}
The absolute magnitude of the cross-session performance loss limits real-world usability on commodity hardware. While error across all conditions was concentrated in Extension$\leftrightarrow$Close confusion, the degradation of Rest recall is particularly problematic. On commodity hardware at Day 29, \cmRestMisclass{} of all Rest windows were misclassified as active gestures (12.8\% Extension, 4.4\% Close). In the absence of temporal smoothing or majority voting, this window-level rate indicates a substantial risk of false activation during intended rest, which is a plausible contributor to device abandonment in myoelectric control. Because windows overlap by 50\%, event-level activation frequency cannot be inferred directly from this rate.

These findings contextualize prior comparisons of consumer and clinical sensors. Pizzolato et al. demonstrated that commodity hardware performs comparably to research-grade equipment~\cite{Pizzolato2017-zd}. Our results are consistent with this finding regarding basic signal quality: the tested reduction to \fsCommodity{}~Hz and 10-bit quantization did not materially drive performance loss. The SeNic benchmark provides further context, characterizing longitudinal and non-ideal operating conditions on a physical Myo armband sampling at up to 200 Hz~\cite{9771219}; that a consumer device operating well below our \fsCommodity{}~Hz commodity condition remains usable is broadly consistent with our finding that the tested reduction from \fsClinical{} to \fsCommodity{}~Hz was not a binding constraint in this task, although the device, task, preprocessing, and dataset all differ and no direct extrapolation to 200~Hz is implied. However, both of these evaluations retain eight or more channels, and both treat the acquisition device as an aggregate variable. Our longitudinal results show that when spatial coverage is reduced to that of a minimal two-channel device, the system becomes substantially more fragile to cross-session temporal drift.

% ============================================================
% LIMITATIONS
% ============================================================
\section{Limitations}

The classification task comprises three classes (wrist extension, hand close, and rest), which is narrower than the seven to seventeen gesture sets common in the sEMG literature. A three-class task is more separable than a larger one, so the absolute accuracies reported here are not directly comparable to multi-gesture benchmarks, and the magnitude of the spatial penalty may differ when more classes must be distinguished from the same two channels. The paired comparisons remain internally valid for this three-class task, since all six conditions share the identical task, although greater task complexity could alter the magnitude of the hardware effects. This study evaluates an able-bodied cohort executing a single recording protocol. The temporal degradation observed does not fully capture the distinct physiological dynamics of residual-limb EMG in amputee populations. Furthermore, spatial correlation values justifying the bipolar derivations were extracted from a representative trial rather than a cohort mean, and individual flexor/extensor dominance is anatomically inconsistent across subjects, which introduces variance into the rigid landmark-anchored channel selection.

The 10-bit result also assumes a full-scale voltage range calibrated to the 99.9th percentile of the Session 1 signal. A physical front end with fixed gain, a fixed reference voltage, or additional headroom would occupy fewer effective ADC codes for the same signal and could therefore exhibit greater quantization loss than the simulated condition. The embedded verification establishes numerical equivalence between the offline analysis and the deployed implementation from ADC-code reconstruction onward. It does not validate the analog front end or physical electrode acquisition, since codes were replayed from recorded data over a deliberately paced serial link rather than acquired by the on-board converter. Continuous \fsCommodity{}~Hz real-time scheduling, including the interrupt-driven architecture identified above as necessary, was neither implemented nor validated. Sample-level filter equality was confirmed on a single trial, and the verification as a whole covers one subject's held-out partition. Separately, the per-stage latency and worst-case execution figures were obtained from a profiling sketch representative of the computational workload rather than from the verified firmware, so they characterize the cost of the pipeline rather than timing measured on the implementation that was numerically validated. The simulated hardware degradation pipeline inherently carries specific assumptions. We enforce a conservative signal fidelity penalty: the pooled 8-channel $V_{MAX}$ threshold quantized the clinical data approximately 1.9$\times$ more coarsely than the 2-channel condition. Additionally, the $V_{MAX}$ calibration phase required exposure to within-session test data (Day 1, Trials 6--7), though longitudinal test sets remained strictly isolated. The pipeline also assumes ideal 60 Hz interference rejection natively present in the GRABMyo dataset, whereas physical commodity sensors (e.g., MyoWare) often lack hardware notch filters. The spatial ablation isolates the extremes of the hardware spectrum (two versus eight channels); a continuous sweep evaluating intermediate channel counts was not performed.

The statistical analysis plan was documented after an initial exploratory run rather than pre-registered in advance. All five model architectures are reported regardless of outcome, and post-hoc observations are labeled as such, but the plan does not constitute pre-registration.

Finally, the study explicitly scopes evaluation to classical machine learning utilizing four fundamental time-domain features. It is plausible that convolutional neural networks operating on raw or spectrally transformed signals might partially recover the spatial gap, though such architectures violate the memory and compute ceilings of the targeted ATmega328P microcontroller. Similarly, our \windowMs{} ms window size enforces a strict latency floor that exceeds optimal myoelectric delay bounds; accelerating the physical compute cannot resolve this constraint, requiring future work to evaluate feature stability at shorter window lengths.

% ============================================================
% CONCLUSION
% ============================================================
\section{Conclusion}

The deployment of surface electromyography classifiers on commodity edge devices requires strict hardware compromises. Within the tested hardware ranges, spatial resolution---not the reduction from \fsClinical{} to \fsCommodity{}~Hz or quantization to 10 bits---was the dominant determinant of cross-session viability. Reducing an 8-channel array to a 2-channel commodity layout exacted an 11-percentage-point penalty to baseline accuracy and added 9 percentage points of cross-session degradation over the 28-day interval.

One candidate explanation is that the spatial constraint removes the training-time redundancy that ensemble architectures can otherwise exploit, although the present design cannot fully separate this from differences in model capacity. Consequently, deploying tree-based ensembles on 2-channel hardware yields diminishing returns that fail to justify their microcontroller memory footprints. Feature attribution remained stable across sessions, consistent with shifted inputs being evaluated against fixed decision boundaries, rather than algorithmic reorganization. To achieve robust cross-session myoelectric control on the edge, hardware design should prioritize physical electrode count and interface stability over computational complexity and signal resolution.

% ============================================================
\section*{Declaration of Generative AI and AI-assisted technologies in the writing process} During the preparation of this work the authors used OpenAI ChatGPT, Anthropic Claude, and Google Gemini in order to draft and revise manuscript text, to assist with the development and review of analysis code, and to identify candidate references. All bibliographic entries were subsequently verified against publisher records by the authors. After using these tools, the authors reviewed and edited the content as needed and take full responsibility for the content of the publication.

\bibliographystyle{elsarticle-num}
\bibliography{refs}

\end{document}